% GENERATED FILE. Edit canonical lessons, then run npm run generate:books.
% canonical-source-hash: fnv1a64:2da826d045e53d91
% canonical-lessons: FA-C40-shodan, FA-C40-sakhtan, FA-C40-mahi, FA-C40-berenj, FA-C40-gusht

\chapter{To become, To build, Fish, Rice, Meat}
\label{ch:fa-to-become-to-build-fish-rice-meat}
\hlchaptermodality{40}

\begin{chapteropening}
\textbf{By the end of this chapter:} \emph{I can say to become, to build, to make, a fish, rice and meat.}

Complete the last lesson of chapter 40: I can say to become, to build, to make, a fish, rice and meat.
\end{chapteropening}

\section[\texorpdfstring{\fa{شدن} (shodan)}{شدن (shodan)}]{\fa{شدن} (shodan) — to become}

\label{lesson:FA-C40-shodan}



\begin{quote}
Before the new one: say the Persian for to come back, then the Persian for to take away; to win.
\end{quote}

\subsection*{You'll want to know: \fa{شدن}}
\textbf{\fa{شدن}} — \emph{shodan} — \textquotedblleft{}to become\textquotedblright{}.

\begin{cousinweb}
\emph{shodan} is a native Persian verb, inherited rather than borrowed; it is usually traced to Middle Persian \emph{šudan}, which meant \textquotedblleft{}to go\textquotedblright{}. \textquotedblleft{}To go\textquotedblright{} turning into \textquotedblleft{}to become\textquotedblright{} is a shift English \emph{go} shows too, as in \textquotedblleft{}go bad\textquotedblright{}.
\end{cousinweb}

\begin{grammarlens}[title={What you've built}]
One more word for this chapter.
\end{grammarlens}

\subsection*{Your turn}
\begin{itemize}
\raggedright
  \item \emph{Say it:} \emph{shodan}
  \item \emph{Say it:} \emph{shodan}, once more
  \item \emph{Say it:} say \emph{shodan}
  \item \emph{Recall:} say the Persian for to come back, then the Persian for to take away; to win, then say \emph{shodan} again
\end{itemize}

\begin{tcolorbox}[breakable,colback=teal!4,colframe=teal!35!black,title={Before you move on}]
What does \fa{شدن} mean? (\textquotedblleft{}To become\textquotedblright{}.) Say it once more.
\end{tcolorbox}
\section[\texorpdfstring{\fa{ساختن} (sâkhtan)}{ساختن (sâkhtan)}]{\fa{ساختن} (sâkhtan) — to build, to make}

\label{lesson:FA-C40-sakhtan}



\begin{quote}
Before the new one: say the Persian for to take away; to win, then the Persian for to become.
\end{quote}

\subsection*{You'll want to know: \fa{ساختن}}
\textbf{\fa{ساختن}} — \emph{sâkhtan} — \textquotedblleft{}to build, to make\textquotedblright{}.

\begin{cousinweb}
\emph{sâkhtan} is a native Persian word, inherited rather than borrowed; it is usually traced to Middle Persian \emph{sāxtan}, \textquotedblleft{}to make, to prepare\textquotedblright{}.
\end{cousinweb}

\begin{grammarlens}[title={What you've built}]
One more word for this chapter.
\end{grammarlens}

\subsection*{Your turn}
\begin{itemize}
\raggedright
  \item \emph{Say it:} \emph{sâkhtan}
  \item \emph{Say it:} \emph{sâkhtan}, once more
  \item \emph{Say it:} say \emph{sâkhtan}
  \item \emph{Recall:} say the Persian for to take away; to win, then the Persian for to become, then say \emph{sâkhtan} again
\end{itemize}

\begin{tcolorbox}[breakable,colback=teal!4,colframe=teal!35!black,title={Before you move on}]
What does \fa{ساختن} mean? (\textquotedblleft{}To build, to make\textquotedblright{}.) Say it once more.
\end{tcolorbox}
\section[\texorpdfstring{\fa{ماهی} (mâhi)}{ماهی (mâhi)}]{\fa{ماهی} (mâhi) — a fish}

\label{lesson:FA-C40-mahi}



\begin{quote}
Before the new one: say the Persian for to become, then the Persian for to build.
\end{quote}

\subsection*{You'll want to know: \fa{ماهی}}
\textbf{\fa{ماهی}} — \emph{mâhi} — \textquotedblleft{}a fish\textquotedblright{}.

\begin{cousinweb}
\emph{mâhi} is a native Persian word, inherited rather than borrowed: it continues Middle Persian \emph{māhīg}, \textquotedblleft{}fish\textquotedblright{} (EtymDB). EtymDB takes it back to an Indo-Iranian word for \textquotedblleft{}fish\textquotedblright{}, the same word as Sanskrit \emph{matsya}.
\end{cousinweb}

\begin{grammarlens}[title={What you've built}]
One more word for this chapter.
\end{grammarlens}

\subsection*{Your turn}
\begin{itemize}
\raggedright
  \item \emph{Say it:} \emph{mâhi}
  \item \emph{Say it:} \emph{mâhi}, once more
  \item \emph{Say it:} say \emph{mâhi}
  \item \emph{Recall:} say the Persian for to become, then the Persian for to build, then say \emph{mâhi} again
\end{itemize}

\begin{tcolorbox}[breakable,colback=teal!4,colframe=teal!35!black,title={Before you move on}]
What does \fa{ماهی} mean? (\textquotedblleft{}A fish\textquotedblright{}.) Say it once more.
\end{tcolorbox}
\section[\texorpdfstring{\fa{برنج} (berenj)}{برنج (berenj)}]{\fa{برنج} (berenj) — rice}

\label{lesson:FA-C40-berenj}



\begin{quote}
Before the new one: say the Persian for to build, then the Persian for a fish.
\end{quote}

\subsection*{You'll want to know: \fa{برنج}}
\textbf{\fa{برنج}} — \emph{berenj} — \textquotedblleft{}rice\textquotedblright{}.

\begin{cousinweb}
\emph{berenj} is usually traced to Middle Persian \emph{brinj}, \textquotedblleft{}rice\textquotedblright{}, and set beside Sanskrit \emph{vrīhi}, \textquotedblleft{}rice\textquotedblright{}; the deeper origin of the rice word is not settled.
\end{cousinweb}

\begin{grammarlens}[title={What you've built}]
One more word for this chapter.
\end{grammarlens}

\subsection*{Your turn}
\begin{itemize}
\raggedright
  \item \emph{Say it:} \emph{berenj}
  \item \emph{Say it:} \emph{berenj}, once more
  \item \emph{Say it:} say \emph{berenj}
  \item \emph{Recall:} say the Persian for to build, then the Persian for a fish, then say \emph{berenj} again
\end{itemize}

\begin{tcolorbox}[breakable,colback=teal!4,colframe=teal!35!black,title={Before you move on}]
What does \fa{برنج} mean? (\textquotedblleft{}Rice\textquotedblright{}.) Say it once more.
\end{tcolorbox}
\section[\texorpdfstring{\fa{گوشت} (gusht)}{گوشت (gusht)}]{\fa{گوشت} (gusht) — meat}

\label{lesson:FA-C40-gusht}



\begin{quote}
Before the new one: say the Persian for a fish, then the Persian for rice.
\end{quote}

\subsection*{You'll want to know: \fa{گوشت}}
\textbf{\fa{گوشت}} — \emph{gusht} — \textquotedblleft{}meat\textquotedblright{}.

\begin{cousinweb}
\emph{gusht} is a native Persian word, inherited rather than borrowed; it is usually traced to Middle Persian \emph{gōšt}, \textquotedblleft{}meat\textquotedblright{}.
\end{cousinweb}

\begin{grammarlens}[title={What you've built}]
One more word for this chapter.
\end{grammarlens}

\subsection*{Your turn}
\begin{itemize}
\raggedright
  \item \emph{Say it:} \emph{gusht}
  \item \emph{Say it:} \emph{gusht}, once more
  \item \emph{Say it:} say \emph{gusht}
  \item \emph{Recall:} say the Persian for a fish, then the Persian for rice, then say \emph{gusht} again
\end{itemize}

\begin{tcolorbox}[breakable,colback=teal!4,colframe=teal!35!black,title={Before you move on}]
What does \fa{گوشت} mean? (\textquotedblleft{}Meat\textquotedblright{}.) Say it once more.
\end{tcolorbox}
